% Curs 10 - Pagina 3 (6-XII-1994).
% Fragment independent pentru % Curs 10 - Pagina 3 (6-XII-1994).
% Fragment independent pentru % Curs 10 - Pagina 3 (6-XII-1994).
% Fragment independent pentru % Curs 10 - Pagina 3 (6-XII-1994).
% Fragment independent pentru \input{courses/pages/curs10-pagina3-ro}.
% Continuarea demonstratiei incepute pe pagina 2.

\subsection*{Curs 10 - Pagina 3}

\noindent (2) \quad
$N\trianglelefteq G\ \Longleftrightarrow\ axa^{-1}\in N$,
\quad $(\forall)\,a\in G$, $(\forall)\,x\in N$.

\medskip
\noindent Prin definiția lui $N_G(H)$,
\[
axa^{-1}\in H, \qquad
(\forall)\,a\in N_G(H),\quad (\forall)\,x\in H.
\]
\noindent $\Longrightarrow\ H\trianglelefteq N_G(H)$.
Pot considera grupul $N_G(H)/H$.

\bigskip
% Un semn izolat/cancelat preceda linia de separatie din manuscris.
\hrule
\bigskip

\noindent Consider $A$ mulțime arbitrară, $|A|<\infty$.

\medskip
\noindent Fie $z_1,\ldots,z_q\in A$ o transversală a relației de echivalență
asociată acțiunii lui $G$ pe $A$, ``$\sim_G$''.

\[
\left\{\begin{aligned}
&i\neq j\ \Longrightarrow\ z_i\not\sim_G z_j,\\
&(\forall)\,x\in A,\ (\exists)\,i\ \text{a.î.}\ x\sim_G z_i.
\end{aligned}\right.
\]

\noindent În aceste condiții avem:
\[
A=\operatorname{Orb}(z_1)\cup\operatorname{Orb}(z_2)
\cup\cdots\cup\operatorname{Orb}(z_q)
\qquad (\text{reuniune disjunctă}).
\]

\[
\Longrightarrow\quad
|A|=\sum_{i=1}^{q}[G:\operatorname{Stab}_G(z_i)].
\]
\noindent \underline{Ecuația claselor acțiunii lui $G$ pe $A$}.

\bigskip
\noindent În particular, dacă $A=G$, considerând acțiunea prin conjugare
a lui $G$ pe $G$ ($G$ grup finit),
\[
\Longrightarrow\quad
\operatorname{ord}(G)=\sum_{i=1}^{q}[G:C_G(z_i)]
\]
\noindent - ecuația claselor de elemente conjugate.

\bigskip
\noindent Fie $x\in G$ a.î.
\[
C_G(x)=G\quad\Longleftrightarrow\quad
|\operatorname{orb}_G(x)|=1,
\]
\noindent adică $ax=xa$, $(\forall)\,a\in G$.

\bigskip
\noindent Notăm:
\[
Z_G:=\{x\in G\mid ax=xa,\ (\forall)\,a\in G\}.
\]

\noindent $Z_G\trianglelefteq G$.
\quad $Z_G$ se numește \underline{centrul grupului $G$}.

\medskip
\noindent $x,y\in Z_G\ \stackrel{?}{\Longrightarrow}\ xy^{-1}\in Z_G$;
\quad și dacă $x\in Z_G$, $a\in G$,
\[
\stackrel{?}{\Longrightarrow}\ axa^{-1}\in Z_G
\qquad (axa^{-1}=xaa^{-1}=x).
\]

\bigskip
\noindent $Z_G$ este inclus în orice transversală; în particular,
$Z_G\subseteq\{z_1,\ldots,z_q\}=T$.

\[
\operatorname{ord}(G)=\operatorname{ord}(Z_G)
+\sum_{\substack{z_i\in T\\z_i\notin Z_G}}[G:C_G(z_i)].
\]
% Pentru z_i in centru, [G:C_G(z_i)]=1 (adnotare "1" in manuscris).
.
% Continuarea demonstratiei incepute pe pagina 2.

\subsection*{Curs 10 - Pagina 3}

\noindent (2) \quad
$N\trianglelefteq G\ \Longleftrightarrow\ axa^{-1}\in N$,
\quad $(\forall)\,a\in G$, $(\forall)\,x\in N$.

\medskip
\noindent Prin definiția lui $N_G(H)$,
\[
axa^{-1}\in H, \qquad
(\forall)\,a\in N_G(H),\quad (\forall)\,x\in H.
\]
\noindent $\Longrightarrow\ H\trianglelefteq N_G(H)$.
Pot considera grupul $N_G(H)/H$.

\bigskip
% Un semn izolat/cancelat preceda linia de separatie din manuscris.
\hrule
\bigskip

\noindent Consider $A$ mulțime arbitrară, $|A|<\infty$.

\medskip
\noindent Fie $z_1,\ldots,z_q\in A$ o transversală a relației de echivalență
asociată acțiunii lui $G$ pe $A$, ``$\sim_G$''.

\[
\left\{\begin{aligned}
&i\neq j\ \Longrightarrow\ z_i\not\sim_G z_j,\\
&(\forall)\,x\in A,\ (\exists)\,i\ \text{a.î.}\ x\sim_G z_i.
\end{aligned}\right.
\]

\noindent În aceste condiții avem:
\[
A=\operatorname{Orb}(z_1)\cup\operatorname{Orb}(z_2)
\cup\cdots\cup\operatorname{Orb}(z_q)
\qquad (\text{reuniune disjunctă}).
\]

\[
\Longrightarrow\quad
|A|=\sum_{i=1}^{q}[G:\operatorname{Stab}_G(z_i)].
\]
\noindent \underline{Ecuația claselor acțiunii lui $G$ pe $A$}.

\bigskip
\noindent În particular, dacă $A=G$, considerând acțiunea prin conjugare
a lui $G$ pe $G$ ($G$ grup finit),
\[
\Longrightarrow\quad
\operatorname{ord}(G)=\sum_{i=1}^{q}[G:C_G(z_i)]
\]
\noindent - ecuația claselor de elemente conjugate.

\bigskip
\noindent Fie $x\in G$ a.î.
\[
C_G(x)=G\quad\Longleftrightarrow\quad
|\operatorname{orb}_G(x)|=1,
\]
\noindent adică $ax=xa$, $(\forall)\,a\in G$.

\bigskip
\noindent Notăm:
\[
Z_G:=\{x\in G\mid ax=xa,\ (\forall)\,a\in G\}.
\]

\noindent $Z_G\trianglelefteq G$.
\quad $Z_G$ se numește \underline{centrul grupului $G$}.

\medskip
\noindent $x,y\in Z_G\ \stackrel{?}{\Longrightarrow}\ xy^{-1}\in Z_G$;
\quad și dacă $x\in Z_G$, $a\in G$,
\[
\stackrel{?}{\Longrightarrow}\ axa^{-1}\in Z_G
\qquad (axa^{-1}=xaa^{-1}=x).
\]

\bigskip
\noindent $Z_G$ este inclus în orice transversală; în particular,
$Z_G\subseteq\{z_1,\ldots,z_q\}=T$.

\[
\operatorname{ord}(G)=\operatorname{ord}(Z_G)
+\sum_{\substack{z_i\in T\\z_i\notin Z_G}}[G:C_G(z_i)].
\]
% Pentru z_i in centru, [G:C_G(z_i)]=1 (adnotare "1" in manuscris).
.
% Continuarea demonstratiei incepute pe pagina 2.

\subsection*{Curs 10 - Pagina 3}

\noindent (2) \quad
$N\trianglelefteq G\ \Longleftrightarrow\ axa^{-1}\in N$,
\quad $(\forall)\,a\in G$, $(\forall)\,x\in N$.

\medskip
\noindent Prin definiția lui $N_G(H)$,
\[
axa^{-1}\in H, \qquad
(\forall)\,a\in N_G(H),\quad (\forall)\,x\in H.
\]
\noindent $\Longrightarrow\ H\trianglelefteq N_G(H)$.
Pot considera grupul $N_G(H)/H$.

\bigskip
% Un semn izolat/cancelat preceda linia de separatie din manuscris.
\hrule
\bigskip

\noindent Consider $A$ mulțime arbitrară, $|A|<\infty$.

\medskip
\noindent Fie $z_1,\ldots,z_q\in A$ o transversală a relației de echivalență
asociată acțiunii lui $G$ pe $A$, ``$\sim_G$''.

\[
\left\{\begin{aligned}
&i\neq j\ \Longrightarrow\ z_i\not\sim_G z_j,\\
&(\forall)\,x\in A,\ (\exists)\,i\ \text{a.î.}\ x\sim_G z_i.
\end{aligned}\right.
\]

\noindent În aceste condiții avem:
\[
A=\operatorname{Orb}(z_1)\cup\operatorname{Orb}(z_2)
\cup\cdots\cup\operatorname{Orb}(z_q)
\qquad (\text{reuniune disjunctă}).
\]

\[
\Longrightarrow\quad
|A|=\sum_{i=1}^{q}[G:\operatorname{Stab}_G(z_i)].
\]
\noindent \underline{Ecuația claselor acțiunii lui $G$ pe $A$}.

\bigskip
\noindent În particular, dacă $A=G$, considerând acțiunea prin conjugare
a lui $G$ pe $G$ ($G$ grup finit),
\[
\Longrightarrow\quad
\operatorname{ord}(G)=\sum_{i=1}^{q}[G:C_G(z_i)]
\]
\noindent - ecuația claselor de elemente conjugate.

\bigskip
\noindent Fie $x\in G$ a.î.
\[
C_G(x)=G\quad\Longleftrightarrow\quad
|\operatorname{orb}_G(x)|=1,
\]
\noindent adică $ax=xa$, $(\forall)\,a\in G$.

\bigskip
\noindent Notăm:
\[
Z_G:=\{x\in G\mid ax=xa,\ (\forall)\,a\in G\}.
\]

\noindent $Z_G\trianglelefteq G$.
\quad $Z_G$ se numește \underline{centrul grupului $G$}.

\medskip
\noindent $x,y\in Z_G\ \stackrel{?}{\Longrightarrow}\ xy^{-1}\in Z_G$;
\quad și dacă $x\in Z_G$, $a\in G$,
\[
\stackrel{?}{\Longrightarrow}\ axa^{-1}\in Z_G
\qquad (axa^{-1}=xaa^{-1}=x).
\]

\bigskip
\noindent $Z_G$ este inclus în orice transversală; în particular,
$Z_G\subseteq\{z_1,\ldots,z_q\}=T$.

\[
\operatorname{ord}(G)=\operatorname{ord}(Z_G)
+\sum_{\substack{z_i\in T\\z_i\notin Z_G}}[G:C_G(z_i)].
\]
% Pentru z_i in centru, [G:C_G(z_i)]=1 (adnotare "1" in manuscris).
.
% Continuarea demonstratiei incepute pe pagina 2.

\subsection*{Curs 10 - Pagina 3}

\noindent (2) \quad
$N\trianglelefteq G\ \Longleftrightarrow\ axa^{-1}\in N$,
\quad $(\forall)\,a\in G$, $(\forall)\,x\in N$.

\medskip
\noindent Prin definiția lui $N_G(H)$,
\[
axa^{-1}\in H, \qquad
(\forall)\,a\in N_G(H),\quad (\forall)\,x\in H.
\]
\noindent $\Longrightarrow\ H\trianglelefteq N_G(H)$.
Pot considera grupul $N_G(H)/H$.

\bigskip
% Un semn izolat/cancelat preceda linia de separatie din manuscris.
\hrule
\bigskip

\noindent Consider $A$ mulțime arbitrară, $|A|<\infty$.

\medskip
\noindent Fie $z_1,\ldots,z_q\in A$ o transversală a relației de echivalență
asociată acțiunii lui $G$ pe $A$, ``$\sim_G$''.

\[
\left\{\begin{aligned}
&i\neq j\ \Longrightarrow\ z_i\not\sim_G z_j,\\
&(\forall)\,x\in A,\ (\exists)\,i\ \text{a.î.}\ x\sim_G z_i.
\end{aligned}\right.
\]

\noindent În aceste condiții avem:
\[
A=\operatorname{Orb}(z_1)\cup\operatorname{Orb}(z_2)
\cup\cdots\cup\operatorname{Orb}(z_q)
\qquad (\text{reuniune disjunctă}).
\]

\[
\Longrightarrow\quad
|A|=\sum_{i=1}^{q}[G:\operatorname{Stab}_G(z_i)].
\]
\noindent \underline{Ecuația claselor acțiunii lui $G$ pe $A$}.

\bigskip
\noindent În particular, dacă $A=G$, considerând acțiunea prin conjugare
a lui $G$ pe $G$ ($G$ grup finit),
\[
\Longrightarrow\quad
\operatorname{ord}(G)=\sum_{i=1}^{q}[G:C_G(z_i)]
\]
\noindent - ecuația claselor de elemente conjugate.

\bigskip
\noindent Fie $x\in G$ a.î.
\[
C_G(x)=G\quad\Longleftrightarrow\quad
|\operatorname{orb}_G(x)|=1,
\]
\noindent adică $ax=xa$, $(\forall)\,a\in G$.

\bigskip
\noindent Notăm:
\[
Z_G:=\{x\in G\mid ax=xa,\ (\forall)\,a\in G\}.
\]

\noindent $Z_G\trianglelefteq G$.
\quad $Z_G$ se numește \underline{centrul grupului $G$}.

\medskip
\noindent $x,y\in Z_G\ \stackrel{?}{\Longrightarrow}\ xy^{-1}\in Z_G$;
\quad și dacă $x\in Z_G$, $a\in G$,
\[
\stackrel{?}{\Longrightarrow}\ axa^{-1}\in Z_G
\qquad (axa^{-1}=xaa^{-1}=x).
\]

\bigskip
\noindent $Z_G$ este inclus în orice transversală; în particular,
$Z_G\subseteq\{z_1,\ldots,z_q\}=T$.

\[
\operatorname{ord}(G)=\operatorname{ord}(Z_G)
+\sum_{\substack{z_i\in T\\z_i\notin Z_G}}[G:C_G(z_i)].
\]
% Pentru z_i in centru, [G:C_G(z_i)]=1 (adnotare "1" in manuscris).
